\section{Experiments}
% debate just claim the results from paper. ontology claims separately. more agent + debate, more agent + zero-shot cot

%We conducted various experiments to applied the more agents mechanism to existing approaches, applying them across a diverse array of benchmarks that span reasoning and generation tasks, and utilizing LLMs of varying scales. These experiments demonstrated that an increased number of agents typically leads to performance improvements. By directly utilizing vanilla more agents mechanism without any additional specialized prompt engineering or using multi-agent framework, the results not only matched but occasionally exceeded the performance of more sophisticated methods. This indicates that the more agents mechanism operates with ubiquitous efficiency. Furthermore, the mechanism exhibited robustness in the face of hyperparameter variations.

Our experiments are structured to assess the sampling and voting method across a broad spectrum of benchmarks, encompassing tasks that challenge both reasoning and generative capacities of LLMs of varying sizes. More precisely, two forms of our method are considered.

\textbf{Vanilla Form.} Our method operates standalone without any prompt engineering or multiple LLM-Agent collaboration frameworks. By querying the LLM multiple times, we gather a set of candidate answers. The final output is determined through majority voting.

\textbf{Integration with Other Methods.} Our method can be integrated with various existing method aimed at enhancing the performance of LLMs. These external methods are treated as modular components, or `black boxes', which can be executed repeatedly to yield a set of responses. Subsequently, majority voting is employed to synthesize these responses into a single consensus-based answer.

Initially, we assess the method in its vanilla form to demonstrate its generalizability. Further, we investigate the method's compatibility by integrating it with established prompt engineering techniques and multiple LLM-Agents collaboration frameworks, demonstrating that our method complements and enhances these methods. Lastly, we conduct a robustness analysis to explore the method's performance stability against variations in hyperparameters of LLMs.

\subsection{Experiment Setup}
\textbf{Tasks and datasets.} We evaluate our method on the following benchmarks from reasoning to generation tasks.
\begin{itemize}
    \item \textbf{Arithmetic Reasoning.} We employ grade-school math problems from the GSM8K dataset \cite{gsm} as one of the test sets, from which we randomly select $100$ cases. Additionally, we utilize the more challenging MATH dataset \cite{math}, which is composed of $7$ subareas and includes $5$ different levels of difficulty, totaling $5,000$ questions. From this dataset, we extract $20$ cases from each level of difficulty, ensuring a comprehensive evaluation across the spectrum of challenges it presents.
    \item \textbf{General Reasoning.} For these tasks, we utilize the MMLU dataset \cite{mmlu}, which contains a wide array of problems across $4$ aspects in $57$ subjects. Additionally, we employ the dataset from the chess state tracking task\footnote{\url{https://github.com/google/BIG-bench/blob/main/bigbench/benchmark_tasks/chess_state_tracking/synthetic_short/task.json}} included in the comprehensive BIG-Bench Benchmark \cite{bbb}. From each dataset, we randomly select $100$ cases for our evaluation.
    \item \textbf{Code Generation.} For the code generation task, we employ the HumanEval benchmark \cite{humaneval}, consisting of 164 human-annotated function-level code completion challenges, each with corresponding unit tests for validation. To implement our method, we compute the BLEU score \cite{bleu} among all pairs of generated candidate answers. The candidate with the highest aggregate BLEU score, indicating the greatest consensus across responses, is then selected as the final output.
\end{itemize}

\textbf{Language models.} We assess our method using language models of four different scales from the Llama2 \cite{llama2} and GPT series. Specifically, we evaluate two versions of Llama2-Chat\footnote{\url{https://github.com/facebookresearch/llama}}, optimized for conversational use cases through alignment techniques, with model sizes of 13B and 70B parameters. Additionally, we include GPT-3.5-turbo and GPT-4 in our evaluation.

\textbf{Methods enhanced by sampling and voting method.} We select various typical methods from two distinct categories and enhanced them with our method: 
\begin{itemize}
    \item \textbf{Prompt Engineering.} We utilize Chain-of-Thought prompting (CoT) \cite{cot}, Zero-Shot Chain-of-Thought prompting (Zero-Shot Cot) \cite{zs-cot}, and more sophisticated methods such as Solo Performance Prompting (SPP) \cite{spp}.  Initially, these methods are applied with a single LLM query. We then increase the number of queries and employed majority voting to determine the most consistent answer as the final response. It is worth mentioning that our experiments with CoT enhanced by our method diverge from CoT-SC \cite{cotsc} in scope. We don't only test on reasoning tasks but also extend our experiments to generative scenarios, such as HumanEval \cite{humaneval}. Additionally, we conduct validations across a wider array of prompt engineering methods, resulting in a more comprehensive work.
    \item \textbf{Multiple LLM-Agents Collaboration.} We select LLM-Debate (debate) \cite{debate} and self-reflection (reflection) \cite{reflexion}. With these methods, we expand the number of participating agents and use majority voting on the collective responses to produce the final answer.
\end{itemize}


% We scaled up three distinct existing methods to demonstrate the performance gains achieved by increasing the number of agents: (i) Chain-of-thought prompting \cite{cot}, hereafter referred to as \textbf{CoT}. Initially, we utilized CoT with a single agent and then scaled up the number of agents by incorporating Self-Consistency \cite{cotsc}, referred to as \textbf{CoT-SC}; (ii) Zero-Shot Chain-of-thought prompting \cite{zs-cot}, referred to as \textbf{Zero-Shot CoT}. We began with single-agent results using Zero-Shot CoT and increased the agent number by applying Self-Consistency referred to as \textbf{Zero-Shot CoT-SC}; (iii) LLM Debate \cite{debate}, referred to as \textbf{Debate}. For this mechanism, we directly scaled up the number of agents.

%We contrasted the more agent mechanism with five unsupervised baselines spanning from ensemble optimization to LLM-%based multi-agent coordination: (i) Chain-of-thought prompting \cite{cot}, referred to as \textbf{CoT}; (ii) Self-%Consistency \cite{cotsc}, referred to as \textbf{CoT-SC}; (iii) Self-Consistency with Zero-Shot CoT \cite{zs-cot}, %referred to as \textbf{Zero-Shot CoT-SC}; (iv) LLM Debate \cite{debate}, referred to as \textbf{Debate}.

\subsection{Main Results}
\label{sec:main_results}
The effectiveness of our method is assessed by averaging the results across $10$ independent runs. During each run, we scale up the ensemble size to $40$.

\textbf{Arithmetic Reasoning} The results in Table \ref{table1} indicate that the vanilla form of our method generally enhances performance across both arithmetic reasoning datasets, with different LLMs experiencing varying performance gains ranging from $12\%$ to $24\%$ on the GSM8K dataset and from $6\%$ to $10\%$ on the MATH dataset. Notably, without any prompt engineering or multiple LLM-Agents cooperation frameworks, the vanilla form of our method yielded the highest accuracy on the GSM8K dataset while using the Llama2-70B model, and on the MATH dataset while using the Llama2-13B model. A surprising finding is that a smaller model combined with our method can surpass a larger model in performance simply by increasing the ensemble size. For instance, the Llama2-13B model with $40$ agents achieved a $59\%$ accuracy on the GSM8K dataset, outperforming the Llama2-70B model with a single agent, which scored $54\%$. When integrated with alternative enhancement methods, our method demonstrated further performance improvements, yielding increases between $10\%$ and $21\%$ on the GSM8K dataset, and between $1\%$ and $15\%$ on the MATH dataset.

\begin{table*}[ht!]
\caption{The results of arithmetic reasoning tasks shown by accuracy. The best performance for each task is shown in bold.}
\label{table1}
\centering
\scalebox{0.85}{
\begin{tabular}{ccllll}
\toprule
                       & \multirow{2}{*}{Method} & \multicolumn{3}{c}{Models}              \\ \cline{3-6} 
                       &                         & Llama2-13B               & Llama2-70B                & GPT-3.5-Turbo         & GPT-4     \\ \midrule
\multirow{8}{*}        & Single            & 0.35                     & 0.54                      & 0.73        & 0.88               \\
                       & \textbf{Ours (Vanilla)}    & 0.59 {\small(+0.24)}     & \textbf{0.74} {\small(+0.20)}      & 0.85 {\small(+0.12)}      \\ \cline{2-6} 
                       & CoT \cite{cot}                     & 0.39                     & 0.57                      & 0.74                  \\
                       & \textbf{CoT + Ours}                  & 0.56 {\small(+0.17)}     & 0.72 {\small(+0.15)}      & 0.84 {\small(+0.10)}  \\ \cline{2-6} 
                       & Zero-Shot CoT \cite{zs-cot}        & 0.40                     & 0.57                      & 0.74                  \\
            GSM8K      & \textbf{Zero-Shot CoT + Ours}        & \textbf{0.61} {\small(+0.21)}     & 0.73 {\small(0.16)}       & \textbf{0.88} {\small(+0.14)}  \\ \cline{2-6}
                       & SPP \cite{spp}          &                      &                       &                   \\
                       & \textbf{SPP + Ours}        &      &        &   \\ \cline{2-6} 
                       & Debate \cite{debate}      & 0.38                     & 0.59                      & 0.83                  \\
                       & \textbf{Debate + Ours}       & 0.48 {\small(+0.10)}     & 0.65 {\small(+0.06)}      & 0.85 {\small(+0.02)}  \\ \cline{2-6} 
                       & Reflection \cite{reflexion}      &                      &                       &                   \\
                       & \textbf{Reflection + Ours}       &      &       &   \\ \midrule
\multirow{8}{*}        & Single            & 0.03                     & 0.05                      & 0.29        & 0.40               \\
                       & \textbf{Ours (Vanilla)}    & \textbf{0.09} {\small(+0.06)}     & 0.11 {\small(+0.06)}      & 0.39 {\small(+0.10)} \\ \cline{2-6} 
                       & CoT \cite{cot}                    & 0.04                     & 0.06                      & 0.28                  \\
                       & \textbf{CoT + Ours}                  & 0.06 {\small(+0.02)}     & \textbf{0.13} {\small(+0.07)}      & \textbf{0.41} {\small(+0.13)}  \\ \cline{2-6} 
                       & Zero-Shot CoT \cite{zs-cot}         & 0.03                     & 0.04                      & 0.25                  \\
            MATH       & \textbf{Zero-Shot CoT + Ours}        & 0.08 {\small(+0.05)}     & 0.10 {\small(+0.06)}      & 0.40 {\small(+0.15)}  \\ \cline{2-6} 
                       & SPP \cite{spp}         &                      &                       &                   \\
                       & \textbf{SPP + Ours}        &      &        &   \\ \cline{2-6} 
                       & Debate \cite{debate}      & 0.05                     & 0.10                      & 0.32                  \\
                       & \textbf{Debate + Ours}       & 0.07 {\small(+0.02)}     & 0.11 {\small(+0.01)}      & 0.36 {\small(+0.04)}  \\ \cline{2-6} 
                        & Reflection \cite{reflexion}      &                      &                       &                   \\
                       & \textbf{Reflection + Ours}       &      &       &   \\ \bottomrule
\end{tabular}
}

\end{table*}

\textbf{General Reasoning} As shown in Table \ref{table2}, a similar trend of improved performance is observed in general reasoning tasks. By using our method in the vanilla form, performance gains range from $1\%$ to $4\%$ in the Chess Move Validity task and from $5\%$ to $11\%$ in the MMLU task. Remarkably, the vanilla form of our method attains the highest accuracy in the MMLU task when utilizing the Llama2-13B and GPT-3.5-Turbo models. Furthermore, integration with other methods generally achieve the performance gains ranging from $1\%$ to $13\%$ in the Chess Move Validity task and from $1\%$ to $11\%$ in the MMLU task. By increasing the ensemble size, small model can outperform large model. For instance, by employing our method, a Llama2-13B model surpass the performance of a single Llama2-70B model both in its vanilla form and when integrated with CoT or Debate method in the Chess Move Validity task.

%there are instances within the MMLU task where neither CoT methods nor the debate mechanism yield performance gains compared to results obtained directly from the LLM without any specialized techniques. For example, when using the Llama2-13B model, CoT methods achieve an accuracy of $42\%$, that is equivalent to the performance without the CoT approach. In contrast, the debate mechanism underperforms with an accuracy of $37\%$ in the same scenario.

%Table \ref{table2} presents the results on general reasoning tasks, where we observe that increasing the number of %agents enhances performance across both datasets. The performance gains range from $1\%$ to $13\%$ for the Chess %Move Validity task, and from $1\%$ to $11\%$ for the MMLU task. Note that the more agents mechanism achieves the %highest accuracy in the MMLU task when employing the Llama2-13B model and GPT-3.5-Turbo model. More surprisingly, %in the MMLU task, there are instances where CoT methods or the debate mechanism do not yield performance gains %compared to results obtained directly from the LLM without any specialized techniques. For instance, when using the %Llama2-13B model, CoT methods achieve an accuracy of $42\%$, that is equivalent to the performance without the CoT %approach. In contrast, the debate mechanism underperforms with an accuracy of $37\%$ in the same scenario.

% \begin{table*}[htbp]
% \centering
% \begin{tabular}{ccllll}
% \midrule
%                                      & \multirow{2}{*}{Method} & \multicolumn{3}{c}{Models}              \\ \cline{3-6} 
%                                      &                         & Llama2-13B                 & Llama2-70B                & GPT-3.5-Turbo         & GPT-4        \\ \midrule
% \multirow{8}{*}{Chess Move Validity} & CoT                     & 0.18                       & 0.10                      & 0.50                          \\
%                                      & CoT-SC                  & \textbf{0.23} {\small(+0.07)}       & 0.11 {\small(+0.01)}      & 0.55 {\small(+0.05)}          \\ \cline{2-6} 
%                                      & Zero-Shot CoT           & 0.15                       & 0.20                      & 0.35                          \\
%                                      & Zero-Shot CoT-SC        & 0.16 {\small(+0.01)}       & \textbf{0.27} {\small(+0.07)}      & 0.48 {\small(+0.13)}          \\ \cline{2-6} 
%                                      & Debate (3 agents)       & 0.18                       & 0.14                      & 0.49                          \\
%                                      & Debate (6 agents)       & 0.19 {\small(+0.01)}       & 0.17 {\small(+0.03)}      & \textbf{0.57} {\small(+0.08)}          \\ \cline{2-6} 
%                                      & Single Agent            & 0.14                       & 0.12                      & 0.51                   & 0.65       \\
%                                      & \textbf{Our Method (Vanilla)}    & 0.18 {\small(+0.04)}       & 0.13 {\small(+0.01)}      & 0.55 {\small(+0.04)}          \\ \midrule
% \multirow{8}{*}{MMLU}                & CoT                     & 0.42                       & 0.56                      & 0.61                          \\
%                                      & CoT-SC                  & 0.43 {\small(+0.01)}       & 0.57 {\small(+0.01)}      & 0.64 {\small(+0.03)}          \\ \cline{2-6} 
%                                      & Zero-Shot CoT           & 0.42                       & 0.54                      & 0.58                          \\
%                                      & Zero-Shot CoT-SC        & 0.48 {\small(+0.06)}       & \textbf{0.65} {\small(+0.11)}      & 0.69 {\small(+0.11)}          \\ \cline{2-6} 
%                                      & Debate (3 agents)       & 0.37                       & 0.56                      & 0.56                          \\
%                                      & Debate (6 agents)       & 0.39 {\small(+0.02)}       & 0.58 {\small(+0.02)}      & 0.67 {\small(+0.11)}          \\ \cline{2-6} 
%                                      & Single Agent            & 0.42                       & 0.55                      & 0.59                    & 0.77     \\
%                                      & \textbf{Our Method (Vanilla)}             & \textbf{0.51} {\small(+0.09)}       & 0.60 {\small(+0.05)}      & \textbf{0.70} {\small(+0.11)}          \\ \midrule
% \end{tabular}
% \caption{The results of general reasoning tasks shown by accuracy. The best performance for each task is shown in bold.}
% \label{table2}
% \end{table*}
\begin{table*}[ht!]
\caption{The results of general reasoning tasks shown by accuracy. The best performance for each task is shown in bold.}
\label{table2}
\centering
\scalebox{0.85}{
\begin{tabular}{ccllll}
\toprule
                                     & \multirow{2}{*}{Method} & \multicolumn{3}{c}{Models}              \\ \cline{3-6} 
                                     &                         & Llama2-13B                 & Llama2-70B                & GPT-3.5-Turbo         & GPT-4        \\ \midrule
\multirow{8}{*}                      & Single            & 0.14                       & 0.12                      & 0.51                   & 0.65       \\
                                     & \textbf{Ours (Vanilla)}    & 0.18 {\small(+0.04)}       & 0.13 {\small(+0.01)}      & 0.55 {\small(+0.04)}          \\  \cline{2-6}
                                     & CoT \cite{cot}                   & 0.18                       & 0.10                      & 0.50                          \\
                                     & \textbf{CoT + Ours} & \textbf{0.23} {\small(+0.07)}       & 0.11 {\small(+0.01)}      & 0.55 {\small(+0.05)}          \\ \cline{2-6} 
                                     & Zero-Shot CoT \cite{zs-cot}          & 0.15                       & 0.20                      & 0.35                          \\
               Chess Move Validity   & \textbf{Zero-Shot CoT + Ours}  & 0.16 {\small(+0.01)}       & \textbf{0.27} {\small(+0.07)}      & 0.48 {\small(+0.13)}          \\ \cline{2-6} 
                                     & SPP \cite{spp}                    &                            &                           &                               \\
                                     & \textbf{SPP + Ours}       &                            &                           &                               \\ \cline{2-6} 
                                     & Debate \cite{debate}      & 0.18                       & 0.14                      & 0.49                          \\
                                     & \textbf{Debate + Ours}     & 0.19 {\small(+0.01)}       & 0.17 {\small(+0.03)}      & \textbf{0.57} {\small(+0.08)}          \\ \cline{2-6} 
                                     & Reflection \cite{reflexion}             &                            &                           &                               \\
                                     & \textbf{Reflection + Ours} &                            &                           &                               \\ \midrule
\multirow{8}{*}                      & Single            & 0.42                       & 0.55                      & 0.59                    & 0.77     \\ 
                                     & \textbf{Ours (Vanilla)}             & \textbf{0.51} {\small(+0.09)}       & 0.60 {\small(+0.05)}      & \textbf{0.70} {\small(+0.11)} \\ \cline{2-6}
                                     & CoT \cite{cot}                    & 0.42                       & 0.56                      & 0.61                          \\
                                     & \textbf{CoT + Ours}          & 0.43 {\small(+0.01)}       & 0.57 {\small(+0.01)}      & 0.64 {\small(+0.03)}          \\ \cline{2-6} 
                                     & Zero-Shot CoT \cite{zs-cot}          & 0.42                       & 0.54                      & 0.58                          \\
                    MMLU             & \textbf{Zero-Shot CoT + Ours}     & 0.48 {\small(+0.06)}       & \textbf{0.65} {\small(+0.11)}      & 0.69 {\small(+0.11)}          \\ \cline{2-6} 
                                     & SPP \cite{spp}                    &                            &                           &                               \\
                                     & \textbf{SPP + Ours}       &                            &                           &                               \\ \cline{2-6} 
                                     & Debate \cite{debate}      & 0.37                       & 0.56                      & 0.56                          \\
                                     & \textbf{Debate + Ours}      & 0.39 {\small(+0.02)}       & 0.58 {\small(+0.02)}      & 0.67 {\small(+0.11)}          \\ \cline{2-6} 
                                     & Reflection \cite{reflexion}             &                            &                           &                               \\
                                     & \textbf{Reflection + Ours} &                            &                           &                               \\ 
\bottomrule
\end{tabular}
}

\end{table*}


\textbf{Code Generation} Consistent with the similar observation in reasoning tasks, Table \ref{table3} illustrates that employing our method in its vanilla form typically results in performance improvements, with gains ranging from $4\%$ to $9\%$. When combined with other techniques, we observe further enhancements, with increases from $2\%$ to $7\%$. However, two notable exceptions are observed when integrated with the debate method with the Llama2-13B and Llama2-70B models, which result in failed cases. Additionally, the GPT-3.5-Turbo model exhibits a performance decline, underperforming even the result of a single agent. This decrease in performance is attributed primarily to the noise generated by referencing answers of other agents during the debate process. The synthesized responses, which incorporate input from multiple agents, disrupt the coherence of the code logic, leading to the observed performance degradation.


\begin{table*}[ht!]
\centering
\caption{The results of code generation shown by pass@1. The best performance for each task is shown in bold. }
\scalebox{0.85}{
\begin{tabular}{ccllll}
\toprule
                           & \multirow{2}{*}{Method} & \multicolumn{3}{c}{Models}                                                                          \\ \cline{3-6} 
                           &                         & \multicolumn{1}{c}{Llama2-13B} & \multicolumn{1}{c}{Llama2-70B} & \multicolumn{1}{c}{GPT-3.5-Turbo} & \multicolumn{1}{c}{GPT-4}\\ \midrule
\multirow{8}{*} & Single            & 0.14                           & 0.24                           & 0.67                               & 0.88            \\
                           & \textbf{Ours (Vanilla)}    & 0.18 {\small(+0.04)}           & \textbf{0.33} {\small(+0.09)}           & 0.73 {\small(+0.06)}           \\ \cline{2-6} 
                           & CoT \cite{cot}                    & 0.13                           & 0.3                            & 0.70                              \\
                           & \textbf{CoT + Ours}                  & 0.20 {\small(+0.07)}           & 0.32 {\small(+0.02)}           & \textbf{0.75} {\small(+0.05)}           \\ \cline{2-6} 
                           & Zero-Shot CoT \cite{zs-cot}          & 0.15                           & 0.23                           & 0.67                              \\
            HumanEval      & \textbf{Zero-Shot CoT + Ours}        & \textbf{0.22} {\small(+0.07)}           & 0.29 {\small(+0.06)}           & 0.74 {\small(+0.07)}           \\ \cline{2-6}
                            & SPP \cite{spp}                    &                            &                           &                               \\
                            & \textbf{SPP + Ours}       &                            &                           &                               \\ \cline{2-6} 
                           & Debate \cite{debate}      & 0                              & 0                              & 0.18                              \\
                           & \textbf{Debate + Ours}      & 0                              & 0                              & 0.24 {\small(+0.06)}           \\ \cline{2-6} 
                           & Reflection \cite{reflexion}             &                            &                           &                               \\
                           & \textbf{Reflection + Ours} &                            &                           &                               \\ \bottomrule
\end{tabular}
}
\label{table3}
\end{table*}

\subsection{Ablation Study}
We conducted an ablation study to assess the impact of changes in various sampling-related hyperparameters on the final outcomes. The experiment was carried out by altering the temperature \textit{T} in temperature sampling \cite{temperature} and the nucleus probability \textit{p} in nucleus sampling \cite{neucler}, using the GPT-3.5-Turbo model over an average of 20 runs. As depicted in Figure \ref{fig:ablation}, the degree of performance improvement and the sensitivity to these hyperparameters exhibit considerable variation across different datasets. Generally, a larger number of agents tends to yield more positive effects; similarly, an increase in sampling randomness is associated with enhanced performance.

%We have established an ablation study to observe the impacts that modifications to various sampling-related hyper-%parameters exert on the eventual outcomes. We conduct the experiment by varying \textit{T} in temperature sampling %\cite{temperature} and \textit{p} in nucleus sampling \cite{neucler}, over GPT-3.5-turbo by averaged 20 runs. %Figure \ref{fig:ablation} shows that the extent of the performance enhancement, as well as the sensitivity to the %hyper-parameters, vary significantly across different datasets. In general terms, the greater the number of agents, %the more beneficial the effect; similarly, increased randomness correlates with improved outcomes.
\begin{figure*}[htbp]
    \centering
    \includegraphics[width=\linewidth]{icml2024/fig/abliation_1.png}
    \caption{Our method significant improves accuracy over various sampling parameters and tasks. In general terms, the greater the ensemble size, the more beneficial the effect; similarly, increased randomness correlates with more performance gains.}
    \label{fig:ablation}
\end{figure*}